% ===================================================================
%  TÁBLA Uimh. 15 — An Margadh. — Na hEarraí Bia.
%  Traduction 2026 d'apres texto/io, controlee sur texto/fr, texto/en
%  et texto/es. Consigne : voir texto/ga/10-tabla-01.tex.
%
%  LE TABLEAU DU MARCHE EST CELUI DE LA NOURRITURE, et c'est ici que
%  la regle du lieu d'apprentissage se verifie le plus nettement de
%  toute la colonne — parce que la nourriture s'apprend entierement a
%  table, c'est-a-dire chez soi. Tout ce qui pousse, nage, court ou se
%  cuit sur cette planche a son mot irlandais, sans exception :
%  « arán » le pain, « im » le beurre, « cáis » le fromage, « feoil »
%  la viande, « iasc » le poisson, « úll » la pomme, « práta » la
%  pomme de terre, « oinniún » l'oignon, « gairleog » l'ail,
%  « oisre » l'huitre, « gliomach » le homard, « scadán » le hareng,
%  « giorria » le lievre, « piasún » le faisan.
%
%  ET CE QUI VIENT DE LOIN GARDE LE NOM DE SON PAYS, ce qui n'est pas
%  un emprunt de plus mais la meme regle vue de l'autre bout :
%  « tae », « seacláid », « caife », « bananaí », « anainn »,
%  « sairdíní », « licéar », « rum », « coinneac », « curaçao ». On
%  n'apprend pas le nom d'un fruit qu'on n'a jamais vu pousser ; il
%  arrive avec le fruit, dans la meme caisse.
%
%  ZERO INVERSION, POUR LA MEME RAISON QU'AUX TABLEAUX 12 ET 14 : ce
%  tableau enumere. Les trois passages construits qui auraient pu
%  couter se rangent seuls — « la butiko di la trip-vendistino (25) »
%  et « la kesto dil fumizita haringi (60) », genitifs postposes comme
%  en ido ; et « la policisto (100) persekutas la legum-kolportistini
%  (99) », ou le VSO met en sujet le renvoi qui vient deja le premier.
% ===================================================================

%%K t15-apar-1 apar
\VUsaut{4.0mm}
\VUtitre{15.0pt}{60}{TÁBLA Uimh. 15}
\VUsaut{3.0mm}

%%K t15-tit-1 sub
\VUpk{11.4pt}{40}{An Margadh. --- Na hEarraí Bia.}
\VUsaut{3.0mm}

%%K t15-01-1 p
1. --- Cruinníonn formhór na gceannaithe a sholáthraíonn dúinn na hearraí is gá dár
gcothú faoin \VUgras{halla} \textsuperscript{(1)} sa \VUgras{mhargadh faoi dhíon} nó
sna sráideanna in aice láimhe.

%%K t15-02-1 p
2. --- Tá an \VUgras{díoltóir muiceola} \textsuperscript{(2)}, a dhíolann
\VUgras{liamhás} \textsuperscript{(3)}, \VUgras{putóg dhubh} \textsuperscript{(4)},
\VUgras{ispín mór} \textsuperscript{(5)}, \VUgras{ispín putóg} \textsuperscript{(6)}
agus \VUgras{blonag} \textsuperscript{(7)}, ina sheasamh in aice le \VUgras{bean an ime
agus na cáise} \textsuperscript{(8)}, a bhfuil a taispeántas lán de \VUgras{cháiseanna
Ollannacha} \textsuperscript{(9)} agus crotal dearg orthu, de \VUgras{cháiseanna
Camembert} \textsuperscript{(10)}, de \VUgras{rothaí Gruyère} \textsuperscript{(11)}
agus de \VUgras{bhloic ime} úra \textsuperscript{(12)}.

%%K t15-03-1 p
3. --- Os a comhair sin tá \VUgras{bean díolta glasraí} \textsuperscript{(13)} ag
tairiscint \VUgras{trátaí} breátha \textsuperscript{(14)} dá custaiméirí, agus
\VUgras{cóilis} \textsuperscript{(15)}, \VUgras{sailéad} \textsuperscript{(16)},
\VUgras{cabáiste cinn} \textsuperscript{(17)}, \VUgras{cúirséid}
\textsuperscript{(19)}, \VUgras{mealbhacáin} \textsuperscript{(18)} agus fiú
\VUgras{beacáin} \textsuperscript{(20)}. Ach tá \VUgras{buachaill beorach}, a stop a
\VUgras{charr seachadta} \textsuperscript{(21)} agus é luchtaithe le \VUgras{bairillí
beorach} \textsuperscript{(22)}, díreach os comhair a taispeántais, ag cur cosc ar a
custaiméirí teacht i ngar di. Tugann bean gharraíodóra ciseán mór \VUgras{bliosán}
\textsuperscript{(23)} chuici.

%%K t15-tit-2 sub
\VUsaut{4.0mm}
\VUpk{10.2pt}{40}{Díol agus Ceannach.}
\VUsaut{3.0mm}

%%K t15-04-1 p
4. --- Tá an-anam sa mhargadh, mar tá uair an \VUgras{cheant} \textsuperscript{(24)}
tagtha. Stadann na \VUgras{mná tí} \textsuperscript{(26)}, a thagann ann chun a gcuid
ceannaigh a dhéanamh, agus iad ag gabháil thar shiopa \VUgras{bhean na bputóg}
\textsuperscript{(25)}, chun \VUgras{putóga goile} nó \VUgras{ceann laoi} a
cheannach.

%%K t15-05-1 p
5. --- Tá \VUgras{cócaire} ramhar \textsuperscript{(27)} ag margáil faoi
\VUgras{thuinnín} an-bhreá le \VUgras{bean an éisc} \textsuperscript{(28)}, a
ardaíonn a praghsanna de réir a tola féin. Fuair sí méid mór \VUgras{eascann, sóil,
breac} agus \VUgras{carbán.} Tá a cuid \VUgras{troisc} agus \VUgras{diúilicíní}
an-úr.

%%K t15-tit-3 sub
\VUsaut{4.0mm}
\VUpk{10.2pt}{40}{Earraí Saora agus Daora.}
\VUsaut{3.0mm}

%%K t15-06-1 p
6. --- Tá an \VUgras{bhean díolta éanlaithe} \textsuperscript{(29)}, atá lonnaithe in
aice léi, an-choinsiasach. Nuair a dhíolann sí \VUgras{turcaí, gé, lacha} nó
\VUgras{sicín} simplí, ní iarrann sí ach an praghas cóir.

%%K t15-07-1 p
7. --- Ag coirnéal na cearnóige, ar an gcéad urlár, tá \VUgras{díoltóir línéadaigh}
\textsuperscript{(30)} lonnaithe le tamall gearr, agus bíonn na ceannaitheoirí mná
cinnte i gcónaí go bhfaighidh siad aige rogha an-bhreá \VUgras{éadaí boird,
naipcíní, naprúin} agus \VUgras{ceirteacha.} Ar an urlár céanna chuir
\VUgras{sceanadóir} \textsuperscript{(31)} a cheardlann ar bun. Cuireann sé faobhar ar
\VUgras{sceana} na mban díolta sa halla agus déanann sé \VUgras{corráin} do na
garraíodóirí as an \VUgras{gcruach} is crua.

%%K t15-tit-4 sub
\VUsaut{4.0mm}
\VUpk{10.2pt}{40}{An Siopa Spíosraí.}
\VUsaut{3.0mm}

%%K t15-08-1 p
8. --- Faoina cheardlann siúd osclaíodh siopa mór bia le tamall gearr. Ní hé an
\VUgras{siopa spíosraí} simplí gan tábhacht \textsuperscript{(32)} a bhíodh ann fadó
é a thuilleadh, siopa nach ndíoladh ach na hearraí coitianta, \VUgras{tae}
\textsuperscript{(33)}, \VUgras{seacláid} \textsuperscript{(34)}, \VUgras{siúcra
bloic} \textsuperscript{(37)} agus \VUgras{gallúnach} \textsuperscript{(38)}. Díoltar
earra ar bith ann, \VUgras{brioscaí briosca}, \VUgras{bia stánaithe}
\textsuperscript{(35)}, \VUgras{licéir} \textsuperscript{(36)}, \VUgras{rum,}
\VUgras{coinneac, kirsch, licéar ainíse} agus \VUgras{curaçao.} Faightear fiadhach
ann: \VUgras{giorria} \textsuperscript{(39)}, \VUgras{smólaigh}
\textsuperscript{(40)}, \VUgras{patraiscí} \textsuperscript{(41)}, \VUgras{gearga}
\textsuperscript{(42)}, \VUgras{lacha fhiáin} \textsuperscript{(43)} nó
\VUgras{piasún} \textsuperscript{(44)}, chomh héasca le torthaí coimhthíocha ar nós
\VUgras{bananaí} \textsuperscript{(46)} agus \VUgras{anainn.}

%%K t15-09-1 p
9. --- Bíonn \VUgras{buachaill} an tsiopa spíosraí \textsuperscript{(45)}, atá
an-lácha, réidh chun an rud a theastaíonn a thabhairt: \VUgras{subh}
\textsuperscript{(47)} chuiríní nó cuinsí, \VUgras{geir} \textsuperscript{(48)},
\VUgras{gircíní} \textsuperscript{(49)} nó \VUgras{sairdíní} \textsuperscript{(50)}
saillte.

%%K t15-09-2 p
Tá sin an-áisiúil do na \VUgras{custaiméirí mná} \textsuperscript{(51)}, a
fhaigheann, in aice leis na hearraí is coitianta, na céadtorthaí is annaimhe agus na
torthaí is blasta: \VUgras{piorraí} súmhara \textsuperscript{(52)},
\VUgras{plumaí} \textsuperscript{(53)}, \VUgras{fíonchaora} \textsuperscript{(54)},
\VUgras{péitseoga} \textsuperscript{(55)}, \VUgras{almóinní} \textsuperscript{(56)}
agus \VUgras{cnónna} \textsuperscript{(57)}.

%%K t15-tit-5 sub
\VUsaut{4.0mm}
\VUpk{10.2pt}{40}{Na Custaiméirí.}
\VUsaut{3.0mm}

%%K t15-10-1 p
10. --- Tá an \VUgras{rís} \textsuperscript{(58)} agus na \VUgras{lintilí}
\textsuperscript{(59)} in aice le ciste na \VUgras{scadán deataithe}
\textsuperscript{(60)}; tá an \VUgras{práta} \textsuperscript{(61)} agus an
\VUgras{phónair} \textsuperscript{(63)} le hais na \VUgras{n-úll}
\textsuperscript{(62)}, na \VUgras{bhfigí} \textsuperscript{(64)}, na \VUgras{bplumaí
triomaithe} \textsuperscript{(65)} agus na \VUgras{gcnónna coill}
\textsuperscript{(66)}. Faightear sa siopa sin \VUgras{uibheacha}
\textsuperscript{(67)} úra agus \VUgras{ológa} \textsuperscript{(68)}; dá bhrí sin
líontar \VUgras{ciseáin} \textsuperscript{(70)} agus \VUgras{líonta lóin}
\textsuperscript{(73)} ar an toirt.

%%K t15-11-1 p
11. --- Tá cáil thuillte bainte amach ag \VUgras{caife} \textsuperscript{(72)} an
chomhlachta den scoth sin i measc na \VUgras{gcailíní freastail}
\textsuperscript{(69)} agus na \VUgras{gcócairí mná}, mar déanann \VUgras{an
seanshiopadóir spíosraí} \textsuperscript{(71)} féin é a róstadh go han-aireach.

%%K t15-tit-6 sub
\VUsaut{4.0mm}
\VUpk{10.2pt}{40}{Na Radharcanna.}
\VUsaut{3.0mm}

%%K t15-12-1 p
12. --- Ní hé gach duine a thagann chun an mhargaidh chun soláthar a dhéanamh. I siúl
pictiúrtha na sráidíne a théann chun an halla feictear na daoine is éagsúla. In aice
le \VUgras{buachaill an bhúistéara} \textsuperscript{(75)} lácha, atá ag brú
\VUgras{bara rotha} \textsuperscript{(74)} roimhe, seo beirt shaighdiúirí ar saoire
agus iad an-bhródúil as a n-éide lonrach a thaispeáint: an \VUgras{hasár}
\textsuperscript{(76)} agus \VUgras{dolman} gorm agus \VUgras{seacó cleiteach} air,
agus an \VUgras{ceannaire zuáivigh,} agus a mhuinchille bolgach agus \VUgras{feas} ar
a cheann.

%%K t15-13-1 p
13. --- Tá an \VUgras{báicéir} \textsuperscript{(80)}, atá ina sheasamh ar leac
dorais an \VUgras{tsiopa aráin} \textsuperscript{(78)}, díreach tar éis taos a
\VUgras{aráin} \textsuperscript{(79)} a fhuint agus na \VUgras{corráin}
\textsuperscript{(81)}, na \VUgras{rollóga} \textsuperscript{(82)} agus na
\VUgras{builíní cruinne} \textsuperscript{(83)} atá san fhuinneog a bhaint amach as an
oigheann.

%%K t15-14-1 p
14. --- Os cionn an tsiopa aráin tá cónaí ar \VUgras{bhean línéadaigh} oilte
\textsuperscript{(84)} a fhostaíonn roinnt \VUgras{bróidnéirí} \textsuperscript{(85)}.
In aice léi tá cónaí ar \VUgras{bhean níocháin is iarnála} \textsuperscript{(86)}, a
chuireann stáirse san \VUgras{éadach} \textsuperscript{(88)} tar éis a nite agus a
thriomaithe, agus a iarnálann é leis an \VUgras{iarann} \textsuperscript{(87)}.

%%K t15-tit-7 sub
\VUsaut{4.0mm}
\VUpk{10.2pt}{40}{Feoil, Iasc agus Glasraí.}
\VUsaut{3.0mm}

%%K t15-15-1 p
15. --- I \VUgras{siopa an bhúistéara} \textsuperscript{(89)}, atá ar urlár na talún,
díoltar gach cineál \VUgras{feola, caoireoil} \textsuperscript{(90)},
\VUgras{mairteoil} \textsuperscript{(94)} nó \VUgras{laofheoil} \textsuperscript{(92)}
mar \VUgras{cheathrúna caorach, stéigeanna, rósta} agus \VUgras{gríscíní.} Gearrann
\VUgras{buachaill an bhúistéara} \textsuperscript{(93)} na feolta sin go léir agus
cuireann sé na \VUgras{cnámha smeara} \textsuperscript{(96)} ar leataobh. Ag doras an
tsiopa tá \VUgras{bean na n-oisrí} \textsuperscript{(97)} a dhíolann \VUgras{oisrí}
\textsuperscript{(98)}. Osclaíonn sí iad má theastaíonn sin.

%%K t15-16-1 p
16. --- Íocann na ceannaithe sin go léir paitinn nó táille áite; dá bhrí sin
tugann \VUgras{an póilín} \textsuperscript{(100)} tóir gan trua ar na \VUgras{mná
bochta díolta glasraí} \textsuperscript{(99)} a chuireann iomaíocht orthu agus iad ag
iompar thart ar a gcairteacha \VUgras{cairéad, raidisí, tornapaí, cainneann} nó
\VUgras{bhurlaí asparagais, piseanna úra, spionáiste, samhadh, biolar} agus
\VUgras{peirsil.}

%%K t15-tit-8 sub
\VUsaut{4.0mm}
\VUpk{10.2pt}{40}{Seachain an póilín!}
\VUsaut{3.0mm}

%%K t15-17-1 p
17. --- Tá a fhios go han-mhaith ag an mbuachaill a dhíolann \VUgras{gairleog}
\textsuperscript{(101)} agus \VUgras{oinniúin,} nach bhfuil cead aigesean ach an
oiread a chuid earraí a dhíol in aice leis an margadh. Teitheann sé, agus é ag rith
faoi bhaol ciseán na \VUgras{bportán,} na \VUgras{cráifisce} agus na \VUgras{ribí
róibéis} a leagan, ciseán \VUgras{an díoltóra} \textsuperscript{(102)}
\VUgras{gliomach spíonach} agus \VUgras{gliomach.}

%%K t15-18-1 p
18. --- Bíonn gach duine ag corraí agus ag callán; ach ní chuireann sin cosc ar ghuth
an \VUgras{ghloineadóra} \textsuperscript{(103)} taistil a chloisteáil go soiléir, ná
ar bhéic \VUgras{fhear an ghuail} \textsuperscript{(104)} a d'fhág a chairtín ar feadh
cúpla nóiméad chun \VUgras{mála guail} \textsuperscript{(105)} a sheachadadh.
